% ============================================================
%  ch05  第二讲 §2.1：Schubert 簇
%  原文行 1255–1466
% ============================================================

\section{第二讲：仿射 Grassmann 簇的更多几何}\label{Lecture II}

本讲研究仿射 Grassmann 簇的一些基本几何。设 $F=k((t))$，$\mO=k[[t]]$。除特别说明外，设 $k$ 代数闭，$\underline G=G\otimes_k \mO$ 是常值约化群概形；在不引起混淆时把 $\Gr_{G}$ 简记作 $\Gr$。关于某些结论对非常值群的推广，见 \cite{PR} 与 \cite[\S\ 6.3]{Z11}；关于 $p$-进群的相应结论，见 \cite{Z14}。

\subsection{Schubert 簇}\label{II:Sch var}

沿用 \S\ref{Intro:group} 的记号。有典范双射
\begin{equation}\label{Cartan decomp2}
G(\mO)\backslash G(F)/G(\mO)\cong \xcoch(T)/W\cong \xcoch(T)^+,
\end{equation}
它可以具体实现如下。本小节余下部分固定嵌入 $T\subset B\subset G$。设 $\mu\in\xcoch(T)$ 是余权，它定义映射 $\mu:F^\times\to T(F)\subset G(F)$。以 $t^\mu=\mu(t)$ 记素元 $t$ 在 $G(F)$ 中的像，并把它看作 $LG$ 的 $k$-点；它依赖于 $t$ 的选取。它在 $\Gr$ 中的像仍记作 $t^\mu$，与素元的选取无关。于是有 Cartan 分解（见 \cite[\S\ 3.3.3]{Ti}）
\begin{equation}\label{Cartan decomp}
G(F)=\bigsqcup_{\mu\in\xcoch(T)^+} G(\mO)t^\mu G(\mO).
\end{equation}
注意双陪集 $G(\mO)t^\mu G(\mO)$ 既不依赖素元 $t$，也不依赖嵌入 $T\subset G$。因此 \eqref{Cartan decomp} 诱导 \eqref{Cartan decomp2}。

现在设 $\mE_1,\mE_2$ 是 $D$ 上两个 $G$-主丛，$\beta: \mE_1|_{D^*}\cong \mE_2|_{D^*}$ 是 $D^*$ 上的同构。选取同构 $\phi_1:\mE_1\cong \mE^0$ 与 $\phi_2:\mE_2\cong \mE^0$，便得到平凡 $G$-主丛的自同构 $\phi_2\beta\phi_1^{-1}\in \Aut(\mE^0|_{D^*})$，从而得到 $G(F)$ 中的一个元素。$\phi_1,\phi_2$ 的不同选取会让该元素左乘、右乘 $G(\mO)$ 中的元素。因此 $\beta$ 给出一个确定的元素
\[\inv(\beta)\in G(\mO)\backslash G(F)/G(\mO)\cong \xcoch(T)^+,\]
有时称为 $\beta$ 的\textbf{相对位置}。

\begin{rmk}\label{non-algebraically closed}
设 $K$ 是包含 $k$ 的域（不必代数闭），$\mE_1, \mE_2$ 是 $D_K$ 上两个 $G$-主丛，$\beta: \mE_1|_{D^*_K}\cong \mE_2|_{D^*_K}$ 是同构。此时仍可如下定义 $\inv(\beta)\in \xcoch(T)^+$。过渡到代数闭包 $\bar K/K$ 之后，上述构造给出
\[\inv(\beta_{\bar K})\in G(\bar K[[t]])\backslash G(\bar K((t)))/G(\bar K[[t]])\cong G(\mO)\backslash G(F)/G(\mO),\]
它与代数闭包 $\bar K$ 的选取无关。因此 $\inv(\beta)=\inv(\beta_{\bar K})$ 有定义。
\end{rmk}

现在设 $\mE_1,\mE_2$ 是 $D_R$ 上两个 $G$-主丛，$\beta: \mE_1|_{D_R^*}\cong \mE_2|_{D_R^*}$。对每个 $x\in \Spec R$，基变换给出
\[(\mE_1|_{D_{k(x)}},\mE_2|_{D_{k(x)}},\beta_{k(x)}: \mE_1|_{D_{k(x)}^*}\simeq \mE_2|_{D_{k(x)}^*}).\]
令 $\inv_x(\beta):=\inv(\beta_{k(x)})$。

\begin{prop}\label{HP}
设 $X=\Spec R$。给定 $\mu\in\xcoch(T)^+$，点集
\[X_{\leq \mu}:=\left\{x\in X\mid \inv_x(\beta)\leq \mu\right\}\]
在 $X$ 中 Zariski 闭。
\end{prop}

\begin{proof}
略述证明。设 $\rho_\chi: G\to \GL(V_\chi)$ 是 $G$ 的最高权 $\chi$ 的有限维最高权表示（例如取 $V_\chi$ 为 $G$ 的最高权 $\chi$ 的 Weyl 模），并令 $V_{\chi,\mE_i}=\mE_i\times^GV_{\chi}$，$i=1,2$，记相应的向量丛。则由定义，所论子集包含于集合
\[X_{V_\chi,\leq \mu}:=\left\{x\in X\mid\rho_\chi(\beta_x)(V_{\chi,\mE_1})\subset t^{-(\chi,\mu)}(V_{\chi,\mE_2})\right\},\]
而后者易见在 $X$ 中闭。事实上，在 $X$ 上局部（Zariski）地可设 $V_{\chi,\mE_i}$（$i=1,2$）都平凡（论证同引理 \ref{et loc triv}）。于是 $\rho_\chi(\beta)$ 由一个方阵给出，其元素为 $a_{ij}=\sum_h a_{ij,h}t^h\in R((t))$。此时 $X_{V_\chi,\leq \mu}$ 由方程 $a_{ij,h}=0$（$h<-(\chi,\mu)$）定义。

若 $G$ 的导群 $G_\der$ 单连通，则
\[X_{\leq \mu}=\bigcap_{V_\chi} X_{V_\chi,\leq \mu},\]
其中交取遍 $G$ 的所有有限维最高权表示。故 $X_{\leq \mu}$ 闭。

一般地，总能找到 $G$ 由环面 $D$ 作的中心扩张 $\tilde G$，使其导群单连通（即所谓的 $z$-扩张，见 \cite[Proposition 3.1]{MS}）。用引理 \ref{et loc triv}，在 $X$ 上 \'etale 局部地把 $\beta$ 提升为 $\tilde G$-主丛的映射 $\tilde \beta$，并把 $\mu$ 提升为 $\tilde G$ 的余权 $\tilde \mu$，便化归到前一种情形。
\end{proof}

定义（球）\textbf{Schubert 簇} $\Gr_{\leq \mu}$ 为 $\Gr_G$ 的闭子集
\[\Gr_{\leq \mu}=\left\{(\mE,\beta)\in \Gr\mid \inv(\beta)\leq \mu\right\},\]
并赋以既约概形结构。由命题 \ref{HP}，
\[\Gr_\mu:=\left\{(\mE,\beta)\in \Gr\mid \inv(\beta)=\mu\right\}=\Gr_{\leq \mu}\setminus \bigcup_{\la<\mu}\Gr_{\leq \la}\]
是 $\Gr_{\leq \mu}$ 的开子集，称为 \textbf{Schubert 胞腔}。

\begin{prop}\label{geom of Sch}~
\begin{enumerate}
\item $\Gr_{\mu}$ 是单个 $L^+G$-轨道，且是维数为 $(2\rho,\mu)$ 的光滑拟射影簇。
\item $\Gr_{\leq \mu}$ 是 $\Gr_\mu$ 的 Zariski 闭包，从而是射影簇。
\end{enumerate}
\end{prop}

\begin{proof}
对 (1)，注意 $t^\mu$ 在 $L^+G$ 作用下的稳定化子是 $L^+G\cap t^\mu L^+G t^{-\mu}$。于是诱导映射
\[L^+G/(L^+G\cap t^\mu L^+G t^{-\mu}) \to LG/L^+G,\quad g \mapsto gt^\mu\]
是局部闭嵌入。由 Cartan 分解 \eqref{Cartan decomp}，其像恰为 $\Gr_\mu$。因此 $\Gr_\mu$ 光滑，且 $\Gr_\mu$ 在 $t^\mu$ 处的切空间可等同于 $\frakg(\mO)/\frakg(\mO)\cap \Ad_{t^\mu}\frakg(\mO)$。但容易看出
\[\frakg(\mO)/\frakg(\mO)\cap \Ad_{t^\mu}\frakg(\mO)=\bigoplus_{(\al,\mu)\geq 0} \frakg_\al(\mO)/t^{(\al,\mu)}\frakg_\al(\mO),\]
其中和取遍 $G$ 的正根 $\al$，$\frakg_\al$ 是对应的根空间。于是简单计算得 $\dim \Gr_\mu=(2\rho,\mu)$。

对 (2)，注意若 $\la\leq \mu$，则存在正余根 $\al$，使得 $\mu-\al$ 支配且 $\la\leq \mu-\al\leq \mu$（例如 \cite[Lemma 2.3]{R3}）。因此只需说明 $t^{\mu-\al}$ 含于 $\Gr_\mu$ 的 Zariski 闭包。为此只需在 $\Gr_{\leq\mu}$ 中构造一条曲线 $C$，使 $t^{\mu-\al}\in C$ 且 $C-\{t^{\mu-\al}\}\subset \Gr_\mu$。这化归为下述 $\SL_2$ 计算。

首先，对每个整数 $m$，令 $t^{\la_m}:=\begin{pmatrix} t^m & 0\\ 0 &1\end{pmatrix}$，看作 $\on{PGL}_2(F)$ 中的元素，并令
\[K_m=\Ad_{t^{\la_m}}(L^+\SL_2)\subset L\SL_2.\]
注意 $K_m$ 作为 $\mO$-模的李代理由 $\{t^me, h, t^{-m}f\}$ 张成，其中 $\{e,h,f\}$ 是标准的 $\fraks\frakl_2$-三元组。又注意
\[\sigma_m:=\begin{pmatrix}0&-t^m\\t^{-m}&0\end{pmatrix}=\begin{pmatrix} t^m & 0\\ 0 &1\end{pmatrix}\begin{pmatrix}0&-1\\1&0\end{pmatrix}\begin{pmatrix} t^{-m} & 0\\ 0 &1\end{pmatrix}\in K_m.\]
以 $L^{>0}\SL_2$ 记求值映射 $\on{ev}:L^+\SL_2\to \SL_2$ 的核（见例 \ref{hyperspacial and Iwahori2}），并令 $K_m^{(1)}=\Ad_{t^{\la_m}}(L^{>0}\SL_2)$。则 $K_m/K_m^{(1)}\cong \SL_2$。

现在设 $i_\al: \SL_2\to G$ 是对应于 $\al$ 的典范同态。取 $m=(\mu,\al)-1$，考虑轨道
\begin{equation}\label{curve C}
C_{\mu,\al}:=Li_\al(K_{m})t^\mu.
\end{equation}
注意 $K_{m}^{(1)}\subset L^+G\cap t^\mu L^+G t^{-\mu}$，故 $C_{\mu,\al}$ 是 $K_{m}/K_m^{(1)}=\SL_2$ 下的齐性空间。容易看出 $C_{\mu,\al}\cong \bP^1$，且 $(L^+G\cap Li_\al(K_m))t^\mu\cong \bA^1\subset \bP^1$。此外 $\bP^1\setminus\bA^1=i_\al(\sigma_m)t^\mu$。由熟知等式
\[\begin{pmatrix}0&-1\\1&0\end{pmatrix}=\begin{pmatrix}1&-t^{-1}\\0&1\end{pmatrix}\begin{pmatrix}1&0\\ t&1\end{pmatrix}\begin{pmatrix}1&-t^{-1}\\0&1\end{pmatrix}\begin{pmatrix}t^{-1}&0\\0&t\end{pmatrix},\]
可见 $i_{\al}(\sigma_m)t^\mu=(t^{\mu-\al} \bmod L^+G)$。因此 $C_{\mu,\al}$ 就是所要的曲线。
\end{proof}

\begin{rmk}
为简单起见，设 $G_\der$ 单连通且 $\cha\ k=0$。此时可定义 $\Gr$ 的闭子概形 $\Gr'_{\leq \mu}$，其 $R$-点分类那些 $(\mE,\beta)$，使得对每个不可约最高权表示 $V_\chi$ 都有 $\beta(V_{\chi,\mE})\subset t^{-(\chi,\mu)}(V_{\chi}\otimes R[[t]])$。由命题 \ref{HP} 的证明知 $\Gr_{\leq \mu}$ 是 $\Gr'_{\leq \mu}$ 的既约子概形。但 $\Gr'_{\leq \mu}=\Gr_{\leq \mu}$ 是否成立仍是个开问题，与 \cite[Conjecture 2.14, Conjecture 2.20]{KWWY} 密切相关。
\end{rmk}

\begin{ex}\label{momega1}
设 $G=\GL_n$，$\mu=(r,0,\ldots,0)$（在标准等同 $\xcoch(T)\cong \bZ^n$ 下）。用 \S\ref{aff GLn} 中 $\Gr_{\GL_n}$ 的格表述，
\[\Gr_{\leq \mu}(R)=\left\{ \Lambda\subset \Lambda_0:=R[[t]]^n\mid \on{rk} \Lambda_0/\Lambda=r\right\}.\]
现设 $r=n$。以 $\mN_n$ 记 $n\times n$ 幂零矩阵构成的簇。有自然映射
\begin{equation}\label{nil to Gr}
\mN_n\to \Gr_{\leq \mu},\quad A\mapsto \La=(t-A)\La_0.
\end{equation}
它确实有定义，因为 $\det(t-A)=t^n$。引理 \ref{Lus obs} 将证明这是开嵌入，特别地 $\Gr_{\leq \mu}$ 给出 $\mN_n$ 的一个紧化。在拓扑空间层面，这最早由 Lusztig 注意到（见 \cite{Lu0}）。注意给出支配余权 $\la\leq \mu=(n,0,\ldots,0)$ 等价于给出 $n$ 的一个分拆，故每个 $\la\leq \mu$ 给出 $\mN_n$ 中的幂零轨道 $\mO_\la$。请读者验证 \eqref{nil to Gr} 把 $\mO_\la$ 送到 $\Gr_\la$。
\end{ex}

继续一般理论。有映射
\begin{equation}\label{parity map}
p:\xcoch(T)\to \bZ/2,\quad \mu\mapsto (-1)^{(2\rho,\mu)}
\end{equation}
它经过 $\xcoch(T)\to \pi_1(G)\to \bZ/2$ 分解，故由定理 \ref{components} 诱导映射 $p:\pi_0(\Gr_G)\to \bZ/2$。

\begin{cor}\label{parity}
Schubert 胞腔 $\Gr_{\mu}$ 位于偶（相应地，奇）分支，即 $p(\Gr_{\mu})=1$（相应地，$-1$），当且仅当 $\dim \Gr_{\mu}$ 为偶（相应地，奇）。
\end{cor}

后文中以 $(LG)_{\leq \mu}$（相应地，$(LG)_\mu$）记 $\Gr_{\leq \mu}$（相应地，$\Gr_\mu$）在 $LG$ 中的原像。注意 $(LG)_{\leq \mu}$ 与 $(LG)_{\mu}$ 都是概形。

对余权 $\mu$，以 $P_\mu$ 记 $G$ 中对应于 $\mu$ 的抛物子群，即由那些满足 $\langle\al,\mu\rangle\leq 0$ 的根 $\al$ 的根子群 $U_\al$ 生成的群。以 $g\mapsto \bar{g}:=\on{ev}(g)$ 记求值映射（见例 \ref{hyperspacial and Iwahori2}）。则有自然投影
\[p_\mu: \Gr_\mu\cong L^+G/(L^+G\cap t^\mu L^+G t^{-\mu})\to G/P_\mu,\quad (gt^\mu \bmod L^+G)\mapsto (\bar{g}\bmod P_\mu).\]
其纤维同构于仿射空间。另一方面，把 $G$ 看作 $L^+G$ 中「常值环路」子群，还有闭嵌入
\[i_\mu: G/P_\mu\cong G\cdot t^\mu\subset \Gr_\mu.\]
注意 $\Aut^+(D)$ 在 $\Gr_G$ 上的作用保持 $\Gr_\mu$\footnote{但 $\Aut(D)$ 在 $\Gr_G$ 上的作用并不保持 Schubert 簇：若 $g\in \Aut(D)$ 把 $t$ 送到 $a_0+t$（$a_0\in R$ 幂零），则 $gt^\mu\not\in \Gr_{\leq \mu}$。}，且在旋转环面 $\bG_m^{\on{rot}}$ 下的不动点子概形是 $G\cdot t^\mu$。

\begin{ex}
设 $G=\GL_n$，$\La_0=k[[t]]^n$，$\mu=(m_1,\ldots,m_n)\in \bZ^n$ 是支配余权（故 $m_i\geq m_{i+1}$）。则映射 $p_\mu: \Gr_\mu\to G/P_\mu$ 把格 $\Lambda\in \Gr_\mu$ 送到 $V=\La_0/t\La_0\cong k^n$ 上如下递减滤过：
\[\on{Fil}^iV=t^{-m_i}\La\cap \La_0/t^{-m_i}\La\cap t\La_0.\]
反之，给定 $V$ 上的递减滤过 $\on{Fil}^\bullet$，映射 $i_\mu$ 把它送到格
\[\Lambda = \sum\on{Fil}^iV\otimes  t^{m_i} \mO\subset V\otimes_kF\cong F^n.\]
\end{ex}

对较小的 $\mu$，下面给出 $\Gr_{\leq \mu}$ 更具体的描述。

回忆支配余权 $\mu$ 称为\textbf{极小}（minuscule）的，若 $\mu\neq 0$ 且对每个正根 $\al$ 有 $(\al,\mu)\leq 1$。注意极小余权在偏序 $\leq$ 下是 $\xcoch(T)^+$ 的极小元。因此

\begin{lem}\label{minusch}
若 $\mu$ 极小，则 $\Gr_{\leq\mu}=\Gr_\mu\cong G/P_\mu$。
\end{lem}

于是 $\Gr_G$ 中的极小 Schubert 簇恰是那些分类对应于极小余权的抛物子群的部分旗簇。例子包括通常的 Grassmann 簇 $\Gr(r,n)$、射影空间中的光滑二次超曲面等。

回忆支配余权 $\mu$ 称为\textbf{拟极小}的，若 $\mu\neq 0$、$\mu$ 不极小，且对每个正根 $\al$ 有 $(\al,\mu)\leq 2$。为简单起见，设 $G_\der$ 单。此时 $\mu$ 就是短支配余根，记作 $\theta$，它是 $\xcoch(T)^+-\{0\}$ 在偏序 $\leq$ 下的极小元。对应于 $\theta$ 的根是最高根 $\theta^\vee$。以 $\mL_\theta=G\times^{P_\theta} k_{\theta^\vee}$ 记 $G/P_\theta$ 上的极丰沛线丛，其中 $P_\theta$ 通过特征 $\theta^\vee$ 作用在一维空间 $k_{\theta^\vee}$ 上。注意 $0\in\xcoch(T)^+$ 是唯一严格小于 $\theta$ 的元素。因此 $\Gr_{\leq \theta}=\Gr_\theta\sqcup \Gr_0$。

\begin{lem}\label{quasiminusch}
设（为简单起见）$G_\der$ 单。若 $\mu=\theta$ 拟极小，则经由投影 $p_\theta$，$\Gr_\theta$ 同构于线丛 $\mL_\theta$ 的全空间，而 $\Gr_{\leq \theta}$ 同构于用 $\mL_\theta$ 给出的 $G/P_\theta$ 的射影嵌入的射影锥。在原点处爆破 $\widetilde{\Gr}_{\leq \theta}= \bP(\mL_{\theta}\oplus \mO)$ 给出奇点消解。
\end{lem}

\begin{proof}见 \cite[\S\ 7]{NP}。
\end{proof}

\begin{rmk}
事实上，$\widetilde{\Gr}_{\leq \theta}$ 是 $\Gr_{\leq \theta}$ 的一种 Demazure 消解，见 \cite[Lemma 2.12]{Z14}。
\end{rmk}

\begin{ex}
若 $G=\SL_2$，则 $\theta=\al$ 是正余根。此时 $\Gr_{\leq \theta}$ 同构于二次曲线 $\{(x,y,z)\mid x^2=yz\}\subset \bP^2$ 的射影锥。
\end{ex}

给定 $G$ 的 $\mu_1,\mu_2\in \xcoch^+$，可用 $L^+G$-主丛 $(LG)_{\leq \mu_1}\to \Gr_{\leq \mu_1}$ 定义 $\Gr_{\leq \mu_1}$ 与 $\Gr_{\leq \mu_2}$ 的扭曲积。也可以把它定义为
\[\Gr_{\leq \mu_1}\tilde{\times}\Gr_{\leq \mu_2}=\left\{(\mE_1,\mE_2,\beta_1,\beta_2)\in\Gr\tilde\times\Gr\mid \inv(\beta_1)\leq \mu_1,\ \inv(\beta_2)\leq \mu_2\right\},\]
它在 $\Gr\tilde\times\Gr$ 中闭，因而可表示。类似地，若 $\mmu=(\mu_1,\ldots,\mu_n)$ 是 $G$ 的支配余权序列，可定义
\[\Gr_{\mmu}:=\Gr_{\mu_1}\tilde\times\cdots\tilde\times\Gr_{\mu_n}\subset \Gr_{\leq\mmu}:=\Gr_{\leq \mu_1}\tilde\times\cdots\tilde\times\Gr_{\leq \mu_n}\subset\Gr\tilde\times\cdots\tilde\times\Gr.\]
令 $|\mmu|=\sum \mu_i$，则卷积映射 \eqref{conv m} 诱导
\begin{equation}\label{conv m2}
m: \Gr_{\leq \mmu}\to \Gr_{\leq |\mmu|}, \quad (\mE_\bullet,\beta_\bullet)\mapsto (\mE_n,\beta_1\cdots\beta_n).
\end{equation}
这个映射类似于 Demazure 消解（只是 $\Gr_{\leq \mmu}$ 一般不光滑）。如 \S\ref{Lecture V} 所示，这些映射的几何极为丰富。这里给出一个例子。

\begin{ex}\label{conv v.s. springer}
沿用例 \ref{momega1} 的记号。设 $\omega_1=(1,0,\ldots,0)$，有卷积映射 $m:\Gr_{\leq(\omega_1,\ldots,\omega_1)}\to\Gr_{\leq \mu}$。另一方面有经典的 Springer 消解 $p:\tilde \mN_r\to\mN_r$（见 \cite[\S~1]{Yu}）。这些映射纳入下述交换图
\begin{equation}\label{sat and springer}
\begin{CD}
\Gr_{\leq(\omega_1,\ldots,\omega_1)}@<<< \Gr_{\leq (\omega_1,\ldots,\omega_1)}^{\square}@>>>\tilde \mN_r\\
@VmVV@VVV@VVpV\\
\Gr_{\leq \mu}@<\pi<<\Gr_{\leq \mu}^{\square}@>\phi>> \mN_r,
\end{CD}
\end{equation}
其中 $\Gr_{\leq \mu}^{\square}(R)$ 分类三元组 $(\mE,\beta,\epsilon)$，这里 $(\mE,\beta)\in\Gr_{\leq \mu}(R)$，而 $\epsilon: \La_0/\La\cong R^r$ 是 $R$-模同构。注意图中两个方块都是笛卡尔方块。事实上，$\Gr_{\leq(\omega_1,\ldots,\omega_1)}$ 中位于 $\La\in \Gr_{\leq \mu}$ 上方的一个点给出链 $\La=\La_r\subset\La_{r-1}\subset\cdots\subset \La_0=k[[t]]^n$；经由标架 $\epsilon$，这样的链给出 $k^r$ 的完全旗，反之亦然。

现设 $r=n$。则 \eqref{nil to Gr} 实际给出 $\phi$ 的一个截面 $\mN_n\to \Gr_{\leq \mu}^{\square}$（因为 $(t-A)\La_0$ 有典范平凡化）。于是得到下述笛卡尔方块
\begin{equation}\label{sat and springer2}
\begin{CD}
\tilde\mN_n@>>>\Gr_{\leq(\omega_1,\ldots,\omega_1)}\\
@VpVV@VVmV\\
\mN_n@>>>\Gr_{\leq \mu}.
\end{CD}
\end{equation}
换言之，此情形下的卷积映射 $m$ 延拓了 $\GL_n$ 的 Springer 消解。
\end{ex}

继续一般讨论。下述基本定理概括了 Schubert 簇奇点的基本事实。

\begin{thm}\label{singularity}
以 $p$ 记 $k$ 的特征指数。若 $p\nmid |\pi_1(G_\der)|$，则 $\Gr_{\leq \mu}$ 正规、Cohen--Macaulay、Gorenstein，且有有理奇点。
\end{thm}

\begin{proof}
这是个困难定理。除 Gorenstein 性外，所有论断都在 \cite[Theorem 8]{Fa} 中证明，另见 \cite[Theorem 0.3]{PR}；关键一步是证明 $\Gr_{\leq \mu}$ 正规。

$\Gr_{\leq \mu}$ 的 Gorenstein 性由 \cite[Equation (241)]{BD} 推出。严格说来，\emph{同上}的论证只在 $k$ 的特征为零（或对 $G$ 为极好素数）时有效。更多细节及向 $p\nmid |\pi_1(G_\der)|$ 情形的推广见 \cite[Theorem 6.11]{Z10}，相关讨论另见注 \ref{symplectic sing}。
\end{proof}

\begin{rmk}
简略讨论 $\underline G$ 为 $G$ 的 Iwahori 群概形时的平行故事，细节见 \cite{Fa,PR}。记 $I=\underline G(\mO)$。此时 Cartan 分解形如
\[I\backslash G(F)/I\simeq \widetilde W,\]
其中 $\widetilde W\simeq \xcoch(T)\rtimes W$ 称为 \textbf{Iwahori Weyl 群}。它是（拟）Coxeter 群，带有 Bruhat 序 $\leq$ 与长度函数 $\ell:\widetilde W\to\bZ$。对 $w\in \widetilde{W}$，以 $S_w$ 记 $\Gr_{\underline G}=\Fl_G$ 上对应的 $L^+\underline G$-轨道，以 $S_{\leq w}$ 记其闭包。则 $S_w$ 同构于维数 $\ell(w)$ 的仿射空间，且 $S_{\leq w}=\cup_{w'\leq w}S_w$。所有 $S_{\leq w}$ 都正规、Cohen--Macaulay 且有有理奇点，但一般不是 Gorenstein。
\end{rmk}
